\section{A Function-Theoretic Definition of Data}
\label{sec:definition}

\subsection{Motivating Example}

Take the five bytes \texttt{0x41 0x52 0x4C 0x49 0x5A}. Read one way, this
is the ASCII string \texttt{"ARLIZ"}. Read another way, it is five
unsigned 8-bit integers, $\{65, 82, 76, 73, 90\}$. Read a third way, it is
simply the first five bytes of some binary file whose format hasn't been
specified yet. The bytes themselves never change across these three
readings---what changes is the \emph{interpretation function} applied to
a fixed set of positions, and that distinction, between the bytes and
whatever is reading them, is exactly the structure this section sets out
to formalize.

\subsection{Formal Definition}

\begin{definition}[Index Set and Value Set]
\label{def:index-value}
An \emph{index set} $I$ is a set whose elements serve as positions or
addresses. A \emph{value set} $V$ is a set whose elements serve as the
possible contents at a position.
\end{definition}

\begin{definition}[Datum]
\label{def:datum}
A \emph{datum} is an element $v \in V$, drawn from some value set $V$,
designated as the content held at a single index.
\end{definition}

\begin{definition}[Collection of Data]
\label{def:collection}
A \emph{collection of data} $D$ over index set $I$ and value set $V$ is a
total function
\[
  D : I \to V
\]
assigning to every index $i \in I$ exactly one value $D(i) \in V$.
\end{definition}

Three conditions have to hold before $D$ counts as well-defined at all,
and each one corresponds to a specific way a real system can fail.

\paragraph{Condition 1 (Totality).}
$D$ must be defined for every $i \in I$. A structure that leaves some
indices unassigned is, formally, a \emph{partial} function, not a total
one---and any operation built on the assumption of totality will misbehave,
silently, the moment it is handed a structure that does not actually have it.

\paragraph{Condition 2 (Functionality).}
For every $i \in I$, $D(i)$ has to name exactly one value in $V$, not a
choice among several. A mapping that assigns two candidate values to the
same index---two fields both claiming the same byte offset under
different schema versions is the usual way this happens in
practice---is not a function in the first place, and therefore is not one
well-defined collection of data, no matter how the rest of the system
treats it.

\paragraph{Condition 3 (Codomain Consistency).}
Every value $D(i)$ has to belong to the \emph{same} $V$, for every $i \in
I$, and $V$ has to be fixed before $D$ is constructed---not discovered
afterward, by whatever happens to read the bytes back.

\begin{definition}[Functional Equality]
\label{def:equality}
Two collections of data $D_1 : I_1 \to V_1$ and $D_2 : I_2 \to V_2$ are
\emph{equal}, written $D_1 = D_2$, if and only if $I_1 = I_2$, $V_1 = V_2$,
and $D_1(i) = D_2(i)$ for every $i \in I_1$.
\end{definition}

This is a stricter condition than matching byte representations, and the
difference is not academic. Two structures can be byte-for-byte identical
under encoding while differing in $V$---the ARLIZ example above is
already an instance of this---and Definition~\ref{def:equality} correctly
calls them unequal, even though a system that hashes raw bytes to check
for duplicates never would.

\subsection{Relation to Sequences, Tuples, and Multisets}

A few structures get used informally as if they were just other words for
``data.'' They are not synonyms; they are special cases, and it is worth
being precise about which constraint each one relaxes or fixes:

\begin{description}
  \item[Sequence / array.]
    The case where $I = \{0, 1, \dots, n-1\}$---a function from a finite,
    contiguous, ordered index set into $V$. An array, in other words, is
    just data over the particular index set that integer addressing
    happens to supply.
  \item[Tuple.]
    The case where $I$ is small and fixed, but $V$ is allowed to vary
    per index, so that $D(i) \in V_i$ rather than a single shared $V$.
    A tuple relaxes Condition~3 \emph{deliberately and explicitly}---which
    is the structural opposite of the unsound, implicit relaxation
    described in Section~\ref{sec:implications}.
  \item[Multiset.]
    The case where $I$ only matters up to permutation. A multiset is a
    quotient of the function-theoretic structure, not a competing
    definition of it.
\end{description}